% An open collection of books, films, and videos that came up around the
% sessions. It is meant to grow. Add your own entry by pull request.
%
% An entry is one title line with a margin note. Copy the block below to
% the end of its section (Reading, Films, or Videos), add the work to
% philosophyofai_references.bib, and fill in the blanks. The "Connects to"
% line names the session themes, or General for a work that spans them. \bibentry prints
% the reference inside the margin note; a \cite there would break.
%
%   \medskip
%   \noindent\textbf{Title}\index{Title} by Author or Director (year).%
%   \entrynote{\likes{0}
%   \newthought{First words} of one line, at most 50 characters,
%   on what the work is about.\\
%   \textit{Connects to: Theme, idea, or thought.}\\
%   \bibentry{bibkey}}
%
% \entrynote sets the margin note like \marginnote, measures how tall it is,
% and leaves that much room in the main column below the title line. Every
% note then stands next to its own entry, and the next entry starts below it.
% The room is glued to the title line, so a page never breaks between them
% and no note runs off the bottom of a page.
% The like counter. \likes{n} opens the margin note with a gray heart and
% n on a line of their own, directly above the \newthought opener. A new
% entry starts at \likes{0}. To like an entry, raise its n by one in a
% pull request. \tuftebreak keeps the gap that \newthought would open above
% the note; the two skips after the counter add up to zero and stop the
% following \newthought from opening it a second time.
\providecommand{\likes}[1]{%
  \tuftebreak\noindent{\color{gray}\ding{170}\,#1}\par\nobreak
  \vskip-\tufteskipamount\vskip\tufteskipamount}
\makeatletter
\providecommand{\entrynote}[1]{%
  \setbox\@tempboxa\vbox{\hsize\marginparwidth
    \@tufte@marginnote@font\@tufte@marginnote@justification
    \@tufte@margin@par\noindent #1}%
  \edef\@tempa{\the\dimexpr\ht\@tempboxa+\dp\@tempboxa-\baselineskip\relax}%
  \marginnote{#1}%
  \@nobreaktrue\par\@nobreakfalse
  \nobreak\vspace{\@tempa}\penalty\z@}
\makeatother
% Reading, Films, and Videos each open a new page, with a bookmark in the
% PDF outline.
% A thin rule under the heading spans the text and the margin.
\providecommand{\collection}[1]{%
  \clearpage
  \pdfbookmark[1]{#1}{collection.#1}%
  \section*{#1}%
  \vspace{-1.2ex}%
  \noindent\makebox[\linewidth][l]{%
    \rule{\dimexpr\linewidth+\marginparsep+\marginparwidth\relax}{0.4pt}}\par
  \nobreak\vspace{2.5ex}}

\chapter{Further Reading \& Watching} \index{Further reading and watching}
\printchapterversion{ad_FurtherReadingWatching}

\marginnote{\newthought{Adding an entry}

\begin{enumerate}
      \item Add the work to \texttt{philosophyofai\_references.bib}. Films
            and videos go in as \texttt{@misc}, books as \texttt{@book}. Give everything a link if possible.

      \item Copy an existing entry in
            \texttt{chapters/\allowbreak ad\_FurtherReadingWatching.tex}
            into the section it belongs to and set its like counter
            to \texttt{\textbackslash likes\{0\}}.

      \item Write at most 50 characters on what it is about in the marginnotes.
      Name the vibe, idea, thought or theme it connects to.

      \item Open a pull request. See the contribution page at the beginning of these notes for details.
\end{enumerate}

\newthought{Or hand it to a coding agent.} Point it at
\texttt{lecturenotes/\allowbreak notes\_philosophyofai/\allowbreak AGENTS.md}.
It asks you what the work is and why it belongs here, then runs the four
steps.}

\newthought{The sessions keep
spilling over} into books, films, and videos that are not on the reading
list. This is the place where we share them.

\newthought{Extend it.} If a book, a film, or a video changed how you think
about one of the four themes, it belongs here, and you are the one to add
it. Every entry is short. It names the work, who made it and when, says in
one line what it is about, and lists the themes it connects to.

\newthought{If you read or watched one of them,} come back and increase
its like counter {\color{gray}\ding{170}} by one, you guessed it, via a PR.

\collection{Reading}

\noindent\textbf{The Hitchhiker's Guide to the Galaxy}\index{The Hitchhiker's Guide to the Galaxy}
by Douglas Adams\index{Adams, Douglas} (1979).%
\entrynote{\likes{1}
\newthought{Forty-two:} an exact answer to an unknown question.\\
\textit{Connects to: Anthropomorphism, Agency.}\\
\bibentry{adams1979hitchhiker}}

\medskip
\noindent\textbf{Sophie's World}\index{Sophie's World} by Jostein Gaarder\index{Gaarder, Jostein} (1991).%
\entrynote{\likes{1}
\newthought{Who are you?} Philosophy from Thales to Sartre.\\
\textit{Connects to: Anthropomorphism, Accountability, Agency, Embodiment.}\\
\bibentry{gaarder1991sofies}}

\medskip
\noindent\textbf{Dune}\index{Dune} by Frank Herbert\index{Herbert, Frank} (1965).%
\entrynote{\likes{1}
\newthought{Thinking machines banned,} humans do their work.\\
\textit{Connects to: Agency, Embodiment.}\\
\bibentry{herbert1965dune}}

\medskip
\noindent\textbf{How the World Really Works}\index{How the World Really Works} by Vaclav Smil\index{Smil, Vaclav} (2022).%
\entrynote{\likes{1}
\newthought{The world runs on} ammonia, steel, cement, plastic.\\
\textit{Connects to: General.}\\
\bibentry{smil2022world}}

\medskip
\noindent\textbf{L'\'Etranger}\index{Etranger@L'\'Etranger} by Albert Camus\index{Camus, Albert} (1942).%
\entrynote{\likes{1}
\newthought{Death ignored} in freedom, faced in a cell.\\
\textit{Connects to: Embodiment.}\\
\bibentry{camus1942etranger}}

\medskip
\noindent\textbf{Le Mythe de Sisyphe}\index{Mythe de Sisyphe, Le} by Albert Camus\index{Camus, Albert} (1942).%
\entrynote{\likes{1}
\newthought{One must imagine} Sisyphus happy.\\
\textit{Connects to: Embodiment.}\\
\bibentry{camus1942mythe}}

\medskip
\noindent\textbf{Hyperion}\index{Hyperion} by Dan Simmons\index{Simmons, Dan} (1989).%
\entrynote{\likes{1}
\newthought{The AIs seceded} and still run the web.\\
\textit{Connects to: Agency, Accountability.}\\
\bibentry{simmons1989hyperion}}

\collection{Films}

\noindent\textbf{Dune} and \textbf{Dune: Part Two}\index{Dune} by Denis Villeneuve (2021, 2024).%
\entrynote{\likes{1}
\newthought{Paul foresees the holy war.} Is he to blame?\\
\textit{Connects to: Accountability, Agency.}\\
\bibentry{villeneuve2021dune}\\
\bibentry{villeneuve2024dune}}

\medskip
\noindent\textbf{Everything Everywhere All at Once}\index{Everything Everywhere All at Once} by Daniel Kwan and Daniel Scheinert (2022).%
\entrynote{\likes{1}
\newthought{Every universe at once,} and nothing matters.\\
\textit{Connects to: General.}\\
\bibentry{kwan2022everything}}

\medskip
\noindent\textbf{Ex Machina}\index{Ex Machina} by Alex Garland (2015).%
\entrynote{\likes{1}
\newthought{It is a Turing Test} with the curtain removed.\\
\textit{Connects to: Anthropomorphism, Agency, Embodiment.}\\
\bibentry{garland2015exmachina}}

\collection{Videos}

\noindent\textbf{The Coming Singularity}\index{The Coming Singularity} by Ray Kurzweil\index{Kurzweil, Ray}, Big Think (2009).%
\entrynote{\likes{1}
\newthought{Brain decoded by 2029,} singularity by 2045.\\
\textit{Connects to: Embodiment.}\\
\bibentry{kurzweil2009coming}}
